\documentclass{ximera}

\ifdefined\HCode
  \newcommand{\alttext}[1]{\HCode{<span class="sr-only">#1</span>}}
\else
  \newcommand{\alttext}[1]{} % Fallback for PDF view
\fi


% MATH -----------------------------------------------------------
\newcommand{\nats}{\mathbb N}
\newcommand{\ints}{\mathbb Z}
\newcommand{\rats}{\mathbb Q}
\newcommand{\reals}{\mathbb R}
\newcommand{\complex}{\mathbb C}
\newcommand{\powerset}{\mathscr P}

%-------------------------------------------------------------

\pgfplotsset{compat=1.18}
\begin{document}
\begin{abstract}
 \end{abstract}

    \title{Exemplar Examples 20}
    
    \maketitle

\begin{enumerate}

\item Let's show that for any integer $n>1$ the Laplace Transform of $f(t)=t^n$ satisfies $\displaystyle \mathcal{L}(t^n)=\frac{n}{s}\mathcal{L}(t^{n-1})$ for $s>0$, by applying the definition, $\mathcal{L}(f)=\displaystyle \int_0^{\infty}f(t)e^{-st}\,dt$. 

\vfill

\item Let's use the above, together with $\mathcal{L}(t)=\frac{1}{s^2}$ to derive a formula or $\mathcal{L}(t^n)$,


\vfill

\newpage


\item Let's consider the piecewise continuous function $f(t)=\left\{\begin{array}{lr}
       4-t & \mbox{ if }0\leq t<4\\
       \\
       0 &  \mbox{ if }t\geq 4
     \end{array}\right.$.



\begin{figure}[h]
    \centering
\alttext{The graph of a piecewise continuous function.}
\begin{tikzpicture}
% Axis
\begin{axis}[
axis x line=middle,
axis y line=middle,
ylabel=$y$,
xlabel=$t$,
xtick={0,1,2,3,4,5},
ytick={0,1,2,3,4,5},
xmax=6,
ymax=5,
xmin=0,
ymin=0
]
% Plots
\addplot[blue,ultra thick,domain=0:4] {4-x};
\addplot[blue,ultra thick,domain=4:6] {0};
\end{axis}
\end{tikzpicture}
    \caption{Graph of a piecewise continuous function}
\end{figure}

Let's determine the Laplace Transform of $f(t)$ by applying the definition,\\ $\mathcal{L}(f)=\displaystyle \int_0^{\infty}f(t)e^{-st}\,dt$.

\end{enumerate}



\end{document}